\documentclass[a4paper,12pt]{article}
\usepackage[T2A]{fontenc}
\usepackage[utf8]{inputenc}
\usepackage[russian]{babel}
\usepackage{amsmath,amssymb,mathtools}
\usepackage{helvet}
\renewcommand{\familydefault}{\sfdefault}
\usepackage{icomma}
\usepackage[a4paper,top=22mm,bottom=24mm,left=23mm,right=23mm,headheight=25pt]{geometry}
\usepackage{microtype}
\usepackage{needspace}
\usepackage{tikz}
\usetikzlibrary{arrows.meta,positioning,shapes.geometric,calc}
\usepackage{booktabs,tabularx,array}
\usepackage{enumitem}
\usepackage{fancyhdr}
\usepackage{setspace}
\usepackage{xcolor}
\usepackage[hidelinks]{hyperref}

\definecolor{accent}{HTML}{008080}
\definecolor{darkaccent}{HTML}{005B61}
\definecolor{textgray}{HTML}{333333}
\definecolor{muted}{HTML}{636D70}
\setlength{\parindent}{0pt}
\setlength{\parskip}{5.5pt plus 1pt minus 1pt}
\onehalfspacing
\setlist[itemize]{leftmargin=6mm,topsep=2pt,itemsep=2pt,parsep=0pt}
\setlist[enumerate]{leftmargin=8mm,topsep=3pt,itemsep=9pt,parsep=0pt}
\pagestyle{fancy}
\fancyhf{}
\fancyhead[L]{\small\color{muted} Дробь от числа и число по его дроби}
\fancyhead[R]{\small\color{muted} Ксения Васильченко}
\fancyfoot[C]{\small\color{muted}\thepage}
\renewcommand{\headrulewidth}{0.45pt}
\renewcommand{\headrule}{\hbox to\headwidth{\color{accent}\leaders\hrule height\headrulewidth\hfill}}
\newcommand{\sect}[1]{%
  \Needspace{6\baselineskip}%
  \vspace{5pt}{\Large\color{darkaccent}#1\par}%
  \vspace{2pt}{\color{accent}\rule{\linewidth}{0.55pt}\par}\vspace{5pt}}
\newcommand{\minititle}[1]{\par\Needspace{5\baselineskip}\vspace{2pt}{\color{darkaccent}\bfseries #1}\par\vspace{2pt}}
\newcommand{\important}[1]{{\color{darkaccent}\bfseries #1}}
\newcommand{\ex}[1]{\minititle{#1}}
\begin{document}
\hypersetup{pdftitle={Чек-лист. Дробь от числа и число по его дроби},pdfauthor={Лёвин Артём Александрович}}

% Титульный лист
\thispagestyle{empty}
\begin{center}
\vspace*{25mm}
{\Large\color{darkaccent} МАТЕМАТИКА\par}
\vspace{17mm}
{\LARGE\color{textgray} Чек-лист\par}
\vspace{10mm}
{\LARGE\color{accent} Дробь от числа\\[3pt] и число по его дроби\par}
\vspace{16mm}
{\normalsize Математика. 6 класс\par}
\vfill
{\large Ксения Васильченко\par}
\vspace{12mm}
{3 ноября 2026 г.\par}
\vspace{18mm}
{\small Лёвин Артём Александрович\\[5pt] эксклюзивно для Ксении Васильченко\par}
\vspace{21mm}
\begin{tikzpicture}[x=1cm,y=1cm]
\draw[accent,line width=.9pt] (-4.1,0) .. controls (-2.3,.65) and (-1.1,-.48) .. (0,0)
 .. controls (1.5,.48) and (2.7,-.65) .. (4.1,0);
\end{tikzpicture}
\vspace{8mm}
\end{center}
\clearpage

\sect{1. Что нужно уметь}
После повторения темы проверьте, можете ли Вы самостоятельно:
\begin{itemize}[label=$\square$,itemsep=4pt]
\item находить заданную дробь от числа;
\item находить исходное число, если известна его дробь;
\item определять, какую часть целого составляет данное число;
\item переводить условие текстовой задачи на язык вычислений;
\item находить неизвестную сторону прямоугольника по площади и другой стороне, затем вычислять периметр.
\end{itemize}

\sect{2. Три величины, которые нужно различать}
\important{Главная идея:} в задачах на дроби важно сначала определить, где целое, где его часть и какая доля целого рассматривается.

Обозначим:
\begin{itemize}
\item $N$ --- \textbf{целое}, или исходное число;
\item $f$ --- \textbf{доля} целого (обыкновенная или десятичная дробь);
\item $d$ --- \textbf{часть} целого, соответствующая этой доле.
\end{itemize}

Между ними существует одна связь:
\[
\boxed{\ d=N\cdot f\ }
\]
Из неё получаем три правила:
\[
\begin{aligned}
\text{Найти часть:}\quad &d=N\cdot f,\\[5pt]
\text{Найти целое:}\quad &N=d:f\quad(f\ne0),\\[5pt]
\text{Найти долю:}\quad &f=d:N\quad(N\ne0).
\end{aligned}
\]
Для задач о длине, площади, времени, массе и количестве предметов целое обычно положительно. Если $0<f<1$, найденная часть меньше целого.

\sect{3. Как читать условие задачи}
\begin{center}
\small
\renewcommand{\arraystretch}{1.4}
\begin{tabularx}{\linewidth}{@{}>{\raggedright\arraybackslash}p{65mm}>{\raggedright\arraybackslash}X@{}}
\toprule
\textbf{Фраза в условии} & \textbf{Что она означает}\\
\midrule
«Найдите $\frac25$ от числа $50$» & $f=\frac25$, $N=50$; требуется часть $d$.\\
«$\frac25$ числа составляют $20$» & $f=\frac25$, $d=20$; требуется целое $N$.\\
«$14$ соток составляют $0,56$ площади участка» & $d=14$, $f=0,56$; весь участок --- целое $N$.\\
«Какую часть от $50$ составляет $20$?» & $d=20$, $N=50$; требуется доля $f$.\\
\bottomrule
\end{tabularx}
\end{center}
\important{Подсказка:} слова «всего», «весь», «всех» часто указывают на целое. Но решение выбирайте по смыслу полного условия.

\sect{4. Три основных типа задач}

\ex{Пример 1. Найти дробь от числа}
Найдите $\dfrac25$ от числа $50$.
\[
 d=50\cdot\frac25=\frac{50\cdot2}{5}=20.
\]
\important{Ответ:} $20$.

\ex{Пример 2. Найти число по его дроби}
Найдите число, если $\dfrac25$ этого числа равны $20$.

Известна \textbf{часть} $d=20$ и \textbf{доля} $f=\dfrac25$. Ищем \textbf{целое}:
\[
N=20:\frac25=20\cdot\frac52=50.
\]
\important{Проверка:} $50\cdot\dfrac25=20$. Следовательно, ответ верен.

\ex{Пример 3. Определить долю целого}
Из $36$ страниц тетради заполнены $9$. Какую часть тетради они составляют?
\[
f=9:36=\frac9{36}=\frac14.
\]
\important{Ответ:} $\dfrac14$ тетради.

\sect{5. Десятичные дроби в задачах на целое}
Доля может быть записана десятичной дробью. Например,
\[
0,6=\frac6{10}=\frac35=60\%.
\]
Это означает шесть десятых, или шестьдесят процентов целого.

\ex{Пример 4. Площадь всех засеянных полей}
Подсолнечником занято $3600$ га, что составляет $0,6$ площади всех засеянных полей. Найдите общую площадь полей.

Часть $d=3600$ га, доля $f=0,6$, ищем целое $N$:
\[
N=3600:0,6=3600:\frac6{10}
=3600\cdot\frac{10}{6}=6000\text{ га}.
\]
\important{Проверка:} $6000\cdot0,6=3600$ га.

\ex{Пример 5. Продолжительность фильма}
Первая часть фильма длится $1,55$ ч и составляет $0,62$ его полной продолжительности. Сколько длится весь фильм?
\[
N=1,55:0,62=\frac{155}{100}:\frac{62}{100}
=\frac{155}{62}=2,5\text{ ч}.
\]
\important{Ответ:} $2,5$ ч, то есть $2$ ч $30$ мин.

\important{Обратите внимание:} $0,62$ целого --- это $62\%$ целого. Записи $0,62$ и $0,62\%$ обозначают разные доли.

\sect{6. Типичные ошибки}
\begin{itemize}
\item \textbf{Умножать вместо деления.} Когда дана \textbf{часть} и её \textbf{доля}, целое находят делением $d:f$.
\item \textbf{Путать часть с целым.} Число рядом со словами «что составляет $0,6$ всей площади» обозначает часть.
\item \textbf{Неверно переводить проценты.} $0,6=60\%$, а $0,6\%=0,006$.
\item \textbf{Забывать о смешанных дробях.} Перед делением $3\dfrac3{14}$ нужно заменить дробью $\dfrac{45}{14}$.
\end{itemize}

\clearpage
\sect{7. Когда в задаче несколько формул}
Иногда дроби нужны, чтобы сначала найти неизвестную величину, а затем ответить на основной вопрос задачи.

\ex{Пример 6. Площадь и периметр прямоугольника}
Площадь прямоугольника равна $19\dfrac27\text{ см}^2$, одна сторона --- $3\dfrac3{14}\text{ см}$. Найдите периметр.

\minititle{Шаг 1. Запишите нужные формулы}
\[
S=a\cdot b,\qquad P=2(a+b).
\]
В формуле периметра нужны обе стороны, а известна только сторона $b$. Сначала найдём $a$ по площади: $a=S:b$.

\minititle{Шаг 2. Найдите неизвестную сторону}
Переведём смешанные числа в неправильные дроби:
\[
19\frac27=\frac{135}{7},\qquad 3\frac3{14}=\frac{45}{14}.
\]
Тогда
\[
\begin{aligned}
a&=\frac{135}{7}:\frac{45}{14}
=\frac{135}{7}\cdot\frac{14}{45}\\[3pt]
&=\frac{135\cdot2}{45}=6\text{ см}.
\end{aligned}
\]

\minititle{Шаг 3. Вычислите периметр}
\[
\begin{aligned}
P&=2\left(6+3\frac3{14}\right)
=2\left(9\frac3{14}\right)\\[3pt]
&=2\cdot\frac{129}{14}=\frac{129}{7}
=18\frac37\text{ см}.
\end{aligned}
\]
\important{Проверка площади:}
\[
6\cdot3\frac3{14}=6\cdot\frac{45}{14}
=\frac{135}{7}=19\frac27\text{ см}^2.
\]
\important{Ответ:} $18\dfrac37$ см.\par
Важно: найденная сторона $a=6$ см --- промежуточный результат; вопрос задачи относится к периметру.


\clearpage
\thispagestyle{fancy}
\begin{center}
{\Large\color{darkaccent} Алгоритм: как выбрать действие\par}
\vspace{7mm}
\end{center}
\begin{center}
\begin{tikzpicture}[>=Stealth, node distance=17mm,
  arr/.style={-{Stealth[length=3mm]},line width=.8pt,draw=accent},
  start/.style={ellipse,draw=accent,line width=.8pt,align=center,text width=9.4cm,inner sep=7pt,font=\small},
  question/.style={diamond,aspect=2.2,draw=accent,line width=.8pt,align=center,text width=4.6cm,inner sep=4pt,font=\small},
  result/.style={rounded corners=3pt,draw=accent,line width=.8pt,align=center,text width=4.2cm,minimum height=30mm,inner sep=6pt,font=\small}]
\node[start] (read) at (0,0) {В условии найдите три величины:\\целое $N$, часть $d$, долю $f$};
\node[question] (choose) at (0,-3.0) {Какую величину\\нужно найти?};
\node[result] (p) at (-5.35,-7.0) {\textbf{Часть $d$}\\[5pt]Известны $N$ и $f$\\[7pt]$d=N\cdot f$};
\node[result] (n) at (0,-7.0) {\textbf{Целое $N$}\\[5pt]Известны $d$ и $f$\\[7pt]$N=d:f$};
\node[result] (f) at (5.35,-7.0) {\textbf{Доля $f$}\\[5pt]Известны $d$ и $N$\\[7pt]$f=d:N$};
\draw[arr] (read.south) -- (choose.north);
\draw[arr] (choose.west) -| (p.north);
\draw[arr] (choose.south) -- (n.north);
\draw[arr] (choose.east) -| (f.north);
\end{tikzpicture}
\end{center}
\vspace{14mm}
\begin{center}
{\large\color{darkaccent} После вычисления проверьте результат}\par
\vspace{4mm}
Если доля меньше единицы, соответствующая ей часть должна быть меньше целого.\par
Для проверки подставьте найденные величины в равенство $d=N\cdot f$.
\end{center}

\clearpage
\sect{Тренировочный блок — Путь А}
\textbf{10 заданий.} Сначала определите, что известно: целое, доля или часть. Затем выполните вычисления. Решения запишите самостоятельно.
\begin{enumerate}
\item Найдите $\dfrac38$ от числа $64$.
\item Найдите $0,35$ от числа $200$.
\item Найдите число, если $\dfrac37$ этого числа равны $15$.
\item Найдите число, если $0,4$ этого числа равны $18$.
\item Какую часть числа $72$ составляет число $18$? Ответ запишите несократимой дробью.
\item В альбоме $45$ страниц, из них заполнены $12$. Какую часть альбома составляют заполненные страницы?
\item В магазин привезли $84$ кг яблок, что составляет $0,3$ всей партии фруктов. Найдите массу всей партии.
\item Огород занимает $15$ соток, что составляет $0,6$ площади всего участка. Найдите площадь участка.
\item Первая часть фильма длится $1,44$ ч и составляет $0,6$ времени всего фильма. Найдите продолжительность фильма в часах.
\item Площадь прямоугольника равна $15\dfrac34\text{ см}^2$, а одна сторона --- $2\dfrac14\text{ см}$. Найдите периметр прямоугольника.
\end{enumerate}

\clearpage
\sect{Ответы к тренировочному блоку}
\renewcommand{\arraystretch}{1.36}
\begin{center}
\begin{tabular}{@{}c@{\hspace{30mm}}l@{}}
\toprule
\textbf{№} & \textbf{Ответ}\\
\midrule
1 & $24$\\
2 & $70$\\
3 & $35$\\
4 & $45$\\
5 & $\dfrac14$\\
6 & $\dfrac4{15}$\\
7 & $280$ кг\\
8 & $25$ соток\\
9 & $2,4$ ч\\
10 & $18\dfrac12$ см\\
\bottomrule
\end{tabular}
\end{center}
\vspace{12mm}
\minititle{Итоговая памятка}
\renewcommand{\arraystretch}{1.32}
\begin{tabularx}{\linewidth}{@{}>{\raggedright\arraybackslash}p{62mm}X@{}}
\toprule
\textbf{В условии известно} & \textbf{Нужное действие}\\
\midrule
Целое и доля & Умножьте: $d=N\cdot f$.\\
Часть и доля & Разделите: $N=d:f$.\\
Часть и целое & Разделите: $f=d:N$.\\
Смешанное число & Переведите в неправильную дробь.\\
Площадь и сторона прямоугольника & Найдите вторую сторону, затем периметр.\\
\bottomrule
\end{tabularx}
\vfill
\begin{center}
{\small\color{muted}Лёвин Артём Александрович --- эксклюзивно для Ксении Васильченко}
\end{center}
\end{document}
